\documentclass[aspectratio=169,10pt]{beamer}
% Presentación de resultados. Las cifras vienen de resultados.tex,
% que genera scripts/09_presentacion.py a partir de data/processed/.
\usepackage[utf8]{inputenc}
\usepackage[T1]{fontenc}
\usepackage[spanish,es-nodecimaldot,es-noshorthands]{babel}
\usepackage{booktabs}
\usepackage{graphicx}
\usepackage{siunitx}
\usetheme{metropolis}
\setbeamertemplate{navigation symbols}{}
\graphicspath{{../figures/}}

\newcommand{\pendiente}{\textcolor{orange}{[pendiente]}}
% Archivo generado por scripts/09_presentacion.py: no editar a mano
\newcommand{\periodoMeteo}{2025-01-01 a 2026-10-31}
\newcommand{\umbralParacas}{10}
\newcommand{\ventanaParacas}{12--18}
\newcommand{\mesCandidato}{2025-01}
\newcommand{\eventosCandidato}{10}
\newcommand{\criterioCandidato}{fijado en config.yaml (vientos\_paracas.mes\_candidato)}
\newcommand{\fuenteCandidato}{METAR}
\newcommand{\eventosMetar}{101}
\newcommand{\eventosEra}{0}
\newcommand{\mesMaxMetar}{2025-01 (10 días)}
\newcommand{\vientoFuerteObs}{11.0}
\newcommand{\vientoFuerteEra}{6.4}
\newcommand{\tablaValidacion}{velocidad (m/s) & 15984 & 0.20 & 1.96 & 0.59 \\
velocidad tarde 12-18 h (m/s) & 4665 & -0.46 & 2.29 & 0.36 \\
velocidad cuando obs > 10 m/s & 287 & -4.57 & 4.68 & 0.48 \\
dirección (grados) & 12399 & -25.28 & 52.61 & -- \\
temperatura (°C) & 15981 & -0.25 & 1.18 & 0.92 \\}
\newcommand{\seriePrincipal}{UHSLC}
\newcommand{\anioUhslc}{2022}
\newcommand{\qcAprobados}{1 de 16}
\newcommand{\pendienteIoc}{0.73}
\newcommand{\factorF}{0.75}
\newcommand{\claseMarea}{mixta, predominantemente semidiurna}
\newcommand{\ampMdos}{31.1}
\newcommand{\ampSdos}{9.2}
\newcommand{\ampNdos}{8.8}
\newcommand{\ampKuno}{20.0}
\newcommand{\ampOuno}{10.3}
\newcommand{\ampPuno}{6.1}
\newcommand{\rangoSicigia}{0.81}
\newcommand{\rangoCuadratura}{0.44}
\newcommand{\rangoMaximo}{1.41}
\newcommand{\rangoSicigiaObs}{1.32}
\newcommand{\rangoCuadraturaObs}{0.72}
\newcommand{\varExplicada}{93.8}
\newcommand{\rmsResiduo}{7.6}
\newcommand{\seriePrincipalNombre}{UHSLC 2022}
\newcommand{\tablaArmonicos}{M2 & 12.42 & 31.0 $\pm$ 0.1 & 337.5 $\pm$ 0.2 & 382060 \\
K1 & 23.93 & 20.0 $\pm$ 0.1 & 28.2 $\pm$ 0.3 & 94660 \\
O1 & 25.82 & 10.3 $\pm$ 0.1 & 347.3 $\pm$ 0.7 & 22187 \\
S2 & 12.00 & 9.2 $\pm$ 0.1 & 350.2 $\pm$ 0.5 & 35827 \\
N2 & 12.66 & 8.8 $\pm$ 0.1 & 297.8 $\pm$ 0.7 & 21084 \\
P1 & 24.07 & 6.1 $\pm$ 0.1 & 26.4 $\pm$ 1.0 & 10813 \\}
\newcommand{\dAmMdosIoc}{-0.9}
\newcommand{\dgMdosIoc}{4.3}
\newcommand{\areaBahia}{\pendiente}
\newcommand{\volBahia}{\pendiente}
\newcommand{\profMedia}{\pendiente}
\newcommand{\profMax}{\pendiente}
\newcommand{\dAmMdosFes}{\pendiente}
\newcommand{\dgMdosFes}{\pendiente}


% Figura con aviso si aún no existe el PNG
\newcommand{\figura}[2]{%
  \IfFileExists{../figures/#1.png}%
    {\includegraphics[#2]{#1}}%
    {\fbox{\parbox[c][3cm][c]{0.8\linewidth}{\centering\small\textcolor{orange}{Figura pendiente:}\\\texttt{\detokenize{figures/#1.png}}}}}}

\title{Caracterización de la bahía de Paracas}
\subtitle{Meteorología, marea y batimetría}
\author{Curso de Oceanografía}
\institute{Universidad Nacional Mayor de San Marcos}
\date{\today}

\begin{document}

\maketitle

\begin{frame}{Contenido}
  \tableofcontents
\end{frame}

% ===================================================================
\section{Área de estudio y datos}

\begin{frame}{Área de estudio}
  \begin{columns}[T]
    \column{0.55\textwidth}
      \centering\figura{batimetria_mapa}{height=0.8\textheight}
    \column{0.43\textwidth}
      \small
      \begin{itemize}
        \item Dominio: \SIrange{13.60}{13.95}{\degree}S, \SIrange{76.15}{76.50}{\degree}O.
        \item Bahía somera abierta al norte, protegida por la península de Paracas.
        \item Estaciones:
          \begin{itemize}
            \item Mareógrafo de Pisco (IOC \texttt{pisco}, UHSLC 683)
            \item Aeropuerto de Pisco (METAR SPSO)
          \end{itemize}
        \item Periodo meteorológico: \periodoMeteo{} (UTC).
      \end{itemize}
  \end{columns}
\end{frame}

\begin{frame}{Datos y métodos}
  \small
  \begin{tabular}{@{}lll@{}}
    \toprule
    Componente & Fuente & Método \\
    \midrule
    Meteorología & ERA5 horario (Copernicus CDS) & viento, HR, radiación, eventos Paracas \\
    Validación & METAR SPSO (IEM ASOS) & sesgo, RMSE, $r$, rosas de viento \\
    Marea observada & UHSLC 683 (research quality), IOC & análisis armónico con UTide \\
    Marea modelada & FES2022b (pyTMD) & armónicos en el mareógrafo y en las fronteras \\
    Batimetría & GEBCO (15$''$) & mapa, perfiles, área y volumen \\
    \bottomrule
  \end{tabular}
  \vspace{1em}
  \begin{itemize}
    \item Todo en UTC; la hora local (UTC$-$5) solo se usa para ciclos diurnos.
    \item Repositorio reproducible: \texttt{config.yaml} + scripts \texttt{01}--\texttt{09}.
  \end{itemize}
\end{frame}

% ===================================================================
\section{Meteorología}

\begin{frame}{Régimen de vientos (ERA5, boca de la bahía)}
  \centering
  \figura{era5_rosa_viento_anual}{height=0.82\textheight}\hfill
  \figura{era5_rosa_viento_estaciones}{height=0.82\textheight}
\end{frame}

\begin{frame}{Ciclo diurno del viento}
  \begin{columns}[T]
    \column{0.6\textwidth}
      \figura{era5_ciclo_diurno_viento}{width=\linewidth}
    \column{0.38\textwidth}
      \small
      \begin{itemize}
        \item Máximo por la tarde y mínimo de madrugada: brisa marina intensificada por el contraste térmico con el desierto.
        \item Viento del S--SO, paralelo a la costa.
        \item Ventana de los vientos Paracas: \ventanaParacas{}~h locales.
      \end{itemize}
  \end{columns}
\end{frame}

\begin{frame}{Climatología mensual}
  \centering\figura{era5_climatologia_mensual}{height=0.85\textheight}
\end{frame}

\begin{frame}{Temperatura, humedad y radiación}
  \centering\figura{era5_series_T_HR_radiacion}{height=0.85\textheight}
\end{frame}

\begin{frame}{Vientos Paracas y mes candidato}
  \centering\figura{validacion_eventos_paracas_mensual}{width=0.85\textwidth,height=0.52\textheight,keepaspectratio}
  \small
  \begin{itemize}
    \item Evento: viento $>$ \umbralParacas{}~m/s entre las \ventanaParacas{}~h locales.
    \item Mes candidato para la simulación (\fuenteCandidato): \textbf{\mesCandidato} (\eventosCandidato{} eventos).
    \item Enero y marzo de 2025 empatan en SPSO; se fijó enero porque es el único mes con marea IOC válida.
  \end{itemize}
\end{frame}

\begin{frame}{Validación ERA5 frente a SPSO}
  \centering\figura{validacion_dispersion_era5_spso}{width=0.92\textwidth,height=0.42\textheight,keepaspectratio}\\[0.3em]
  \footnotesize
  \begin{tabular}{@{}lrrrr@{}}
    \toprule
    Variable & $n$ & Sesgo & RMSE & $r$ \\
    \midrule
    \tablaValidacion
    \bottomrule
  \end{tabular}\\[0.3em]
  \scriptsize Sesgo = ERA5 $-$ observado (m/s, grados o $^\circ$C).
\end{frame}

\begin{frame}{¿ERA5 subestima los vientos Paracas?}
  \begin{columns}[T]
    \column{0.55\textwidth}
      \centering
      \figura{validacion_ciclo_diurno_era5_spso}{width=\linewidth,height=0.8\textheight,keepaspectratio}
    \column{0.43\textwidth}
      \small
      \begin{itemize}
        \item Días con evento: SPSO = \textbf{\eventosMetar}, ERA5 = \textbf{\eventosEra}.
        \item Cuando SPSO supera \umbralParacas{}~m/s, ERA5 da \vientoFuerteEra{}~m/s frente a \vientoFuerteObs{}~m/s observados.
        \item \textbf{Sí: ERA5 subestima los eventos}. Con celdas de $\sim$30~km suaviza el contraste tierra--mar y la canalización local del viento.
        \item Para forzar AEM3D conviene corregir el viento de ERA5 o usar SPSO.
      \end{itemize}
  \end{columns}
\end{frame}

\begin{frame}{Rosas de viento: observado y ERA5}
  \centering\figura{validacion_rosas_era5_spso}{height=0.85\textheight}
\end{frame}

% ===================================================================
\section{Marea}

\begin{frame}{Datos del mareógrafo y control de calidad}
  \begin{columns}[T]
    \column{0.6\textwidth}
      \figura{marea_qc_ioc_vs_referencia}{width=\linewidth}
    \column{0.38\textwidth}
      \small
      \begin{itemize}
        \item IOC (tiempo real, sin QC) comparado mes a mes con la marea predicha a partir de UHSLC \anioUhslc.
        \item \texttt{pisco}: \qcAprobados{} meses aprobados. Desde agosto de 2025 las amplitudes son el \pendienteIoc{} del valor esperado: sensor mal calibrado o reubicado.
        \item Serie principal: \textbf{\seriePrincipal{} \anioUhslc} (99,8~\% completa).
      \end{itemize}
  \end{columns}
\end{frame}

\begin{frame}{Constituyentes armónicas (\seriePrincipalNombre)}
  \begin{columns}[T]
    \column{0.55\textwidth}
      \footnotesize\setlength{\tabcolsep}{4pt}
      \begin{tabular}{@{}lrrrr@{}}
        \toprule
        Const. & Periodo (h) & $A$ (cm) & $g$ ($^\circ$) & SNR \\
        \midrule
        \tablaArmonicos
        \bottomrule
      \end{tabular}\\[0.5em]
      \scriptsize Intervalos de confianza del 95~\%; fase de Greenwich (UTC).
    \column{0.43\textwidth}
      \small
      \begin{itemize}
        \item Factor de forma: $F = \dfrac{K_1+O_1}{M_2+S_2} = \mathbf{\factorF}$
        \item Marea \textbf{\claseMarea}.
        \item Varianza explicada: \varExplicada~\%; RMS del residuo: \rmsResiduo~cm.
        \item IOC enero 2025 vs UHSLC: $\Delta A_{M_2}$ = \dAmMdosIoc~cm, $\Delta g_{M_2}$ = \dgMdosIoc$^\circ$.
      \end{itemize}
  \end{columns}
\end{frame}

\begin{frame}{Amplitudes y rangos de marea}
  \begin{columns}[T]
    \column{0.52\textwidth}
      \figura{marea_amplitudes_constituyentes}{width=\linewidth}
    \column{0.46\textwidth}
      \footnotesize\setlength{\tabcolsep}{4pt}
      \begin{tabular}{@{}lrr@{}}
        \toprule
        Rango (m) & Teórico & Observado \\
        \midrule
        Sicigia $2(M_2+S_2)$ & \rangoSicigia & \rangoSicigiaObs \\
        Cuadratura $2(M_2-S_2)$ & \rangoCuadratura & \rangoCuadraturaObs \\
        $2(M_2+S_2+K_1+O_1)$ & \rangoMaximo & -- \\
        \bottomrule
      \end{tabular}\\[0.5em]
      El rango observado supera $2(M_2+S_2)$ porque las diurnas $K_1$ y $O_1$ se suman (marea mixta).
  \end{columns}
\end{frame}

\begin{frame}{Observado, predicho y residuo}
  \centering\figura{marea_observado_predicho_uhslc}{height=0.85\textheight}
\end{frame}

\begin{frame}{Espectro del residuo}
  \begin{columns}[T]
    \column{0.6\textwidth}
      \figura{marea_espectro_residuo}{width=\linewidth}
    \column{0.38\textwidth}
      \small
      \begin{itemize}
        \item El residuo concentra la energía en bajas frecuencias (forzamiento meteorológico, ondas atrapadas a la costa).
        \item Los picos remanentes cerca de 1 y 2~cpd indican variabilidad no astronómica o constituyentes no resueltas.
      \end{itemize}
  \end{columns}
\end{frame}

\begin{frame}{Modelo global FES2022}
  \begin{columns}[T]
    \column{0.5\textwidth}
      \figura{marea_fes_vs_utide}{width=\linewidth}
    \column{0.48\textwidth}
      \figura{marea_armonicos_fronteras}{width=\linewidth}\\
      \small $M_2$: FES2022 $-$ observado: $\Delta A$ = \dAmMdosFes~cm, $\Delta g$ = \dgMdosFes$^\circ$.
  \end{columns}
\end{frame}

\begin{frame}{Forzamiento preparado para AEM3D}
  \begin{columns}[T]
    \column{0.55\textwidth}
      \figura{marea_prediccion_mes_candidato}{width=\linewidth}
    \column{0.43\textwidth}
      \small
      \begin{itemize}
        \item Mes candidato: \textbf{\mesCandidato} (definido con \fuenteCandidato).
        \item Marea de FES2022 cada 15~min en las fronteras norte y oeste (CSV).
        \item Meteorología horaria de ERA5 en la boca (CSV).
        \item \textbf{No se corrió AEM3D}: el formato de entrada se definirá con la plantilla del curso.
      \end{itemize}
  \end{columns}
\end{frame}

% ===================================================================
\section{Batimetría}

\begin{frame}{Perfiles y morfometría}
  \begin{columns}[T]
    \column{0.55\textwidth}
      \figura{batimetria_perfiles}{height=0.8\textheight}
    \column{0.43\textwidth}
      \small
      \begin{tabular}{@{}lr@{}}
        \toprule
        Área de agua & \areaBahia~km$^2$ \\
        Volumen & \volBahia~hm$^3$ \\
        Profundidad media & \profMedia~m \\
        Profundidad máxima & \profMax~m \\
        \bottomrule
      \end{tabular}\\[0.5em]
      Dentro del polígono editable \texttt{data/poligono\_bahia.geojson}. GEBCO es grueso en la bahía; está lista la fusión con sondajes de la carta DHN.
  \end{columns}
\end{frame}

% ===================================================================
\section{Conclusiones}

\begin{frame}{Conclusiones y limitaciones}
  \small
  \begin{itemize}
    \item Viento dominante del S--SO con un fuerte ciclo diurno; los vientos Paracas se concentran por la tarde.
    \item ERA5 reproduce la dirección y el ciclo diurno, pero \textbf{subestima la intensidad} de los eventos (\eventosEra{} frente a \eventosMetar{} días con evento).
    \item Marea \claseMarea{} ($F$ = \factorF), con un rango de sicigia observado de $\sim$\rangoSicigiaObs~m.
    \item El mareógrafo IOC de 2025--26 registra amplitudes atenuadas: para el análisis se usó UHSLC \anioUhslc.
    \item Pendiente: batimetría de detalle (carta DHN), corrección del viento de ERA5 y exportación al formato de AEM3D.
  \end{itemize}
\end{frame}

\end{document}
